\documentclass[../Probability.tex]{subfiles}

\begin{document}

\section{Introduction}
A sample space, $\Omega$, is the set of all possible outcomes of some experiment. \\
An event is a subset of the sample space.

\subsection{Examples}
\subsubsection*{Discrete Sample Space}
Experiment: Roll a die until a '1' appears. \\
Sample space $\Omega = \{(1), (2, 1), (3, 1), \ldots, (2, 2, 1), \ldots, (2, 3, 4, 5, 6, 1, \ldots)\}$. \\
Event that the sum of values is 5: $\{(4, 1), (2, 2, 1)\}$ \\
\textit{Tuples }() \textit{consider the order of elements in it. Sets \{\} do not.}

\subsubsection*{Continuous Sample Space}
Experiment: Output any real number in $[0, 1)$. \\
Sample space $\Omega = [0, 1)$. \\
Example of Events: $\emptyset$, [0, 1), [0.1, 0.3], [0.1, 0.3] $\cup$ (0.5, 0.7) \\
\textit{For continuous sample space, we should define events in  terms of (operations of) intervals.}

\subsection{Probability Distribution}
\paragraph{Definition} Let $\Omega$ be the sample space and $\mathcal{E}$ be the set of all events (where an event is a subset of $\Omega$).
A probability distribution/measure is a function $P: \mathcal{E} \rightarrow [0,1]$,
if it satisfies the following axioms:
\begin{enumerate}
    \item $0 \leq P(A) \leq 1, \quad \forall A \in \mathcal{E}$
    \item $P(\Omega) = 1$
    \item $A_1, A_2, \ldots \text{ are disjoint events in }\mathcal{E} \Rightarrow P(\cup_{i=1}^{\infty} A_i) = \sum_{i=1}^{\infty} P(A_i)$
\end{enumerate}

\paragraph{Properties}
\begin{enumerate}
    \item $P(\emptyset) = 0$
    \item $A \cap B = \emptyset \Rightarrow P(A \cup B) = P(A) + P(B)$
    \item $P(A^C) = 1 - P(A)$
    \item $A \subset B \Rightarrow P(A) \leq P(B)$
    \item $P(A \cup B) = P(A) + P(B) - P(A \cap B)$
\end{enumerate}

\subsection{Independent vs Disjoint events}
An independent event does not affect the probability of occurence of another independent event. They are related by
\[P(A \cap B) = P(A)P(B)\]
Disjoint events can never occur at the same time.
If A and B are disjoint events, then
\[A \cap B = \emptyset \Rightarrow P(A \cap B) = 0\]
In general, independent events are not disjoint and disjoint events are not independent.
The only exception is the trivial case, when one of the events is impossible i.e. $P(A)=0$, hence satisfying both equations above.

\end{document}